\section{Tire Model: Pacejka Magic Formula}

This section presents the Pacejka magic formula, the industry-standard tire model for vehicle dynamics simulation.

\cite{pacejka2005tire} is the authoritative reference. VerletDrift implements a simplified E=0 variant for real-time computation, neglecting pressure effects and asymmetric properties.

\subsection{Pacejka Magic Formula (E=0 Variant)}

\begin{equation}
  F = \frac{N \cdot \mu}{sin(C \cdot \arctan(B))} \cdot \sin(C \cdot \arctan(B \cdot s_{\text{norm}}))
  \label{eq:pacejka-force}
\end{equation}

\textbf{Physical Meaning}: The magic formula converts slip (angle or ratio) into a tire force that peaks at a characteristic slip value then falls off. This nonlinear behavior enables both grip and sliding dynamics.

\textbf{Source Code}: \coderef{physics.js}{506--517}

\begin{lstlisting}[caption={Pacejka magic formula (E=0)}]
const normalisedSlip = slipValue / peakSlipValue;
const peakNorm = Math.sin(C * Math.atan(B));
return (normalLoad * frictionCoeff / peakNorm) *
       Math.sin(C * Math.atan(B * normalisedSlip));
\end{lstlisting}

\textbf{Parameters}:
\begin{itemize}
  \item $N$: Normal load (tire load in Newtons)
  \item $\mu$: Friction coefficient (default: 1.2)
  \item $s_{\text{norm}} = slip / slip_{\text{peak}}$: Normalized slip
  \item $B = \mathrm{pacejkaB}$: Stiffness factor (default: 10)
  \item $C = \mathrm{pacejkaC}$: Shape factor (default: 1.3)
  \item $peakNorm = \sin(C \cdot \arctan(B))$: Normalization to ensure $F_{\text{peak}} = N \cdot \mu$
\end{itemize}

\subsection{Slip Angle}

\begin{equation}
  \alpha = \arctan\left(\frac{v_{\text{lateral}}}{max(|v_{\text{longitudinal}}|, \epsilon)}\right)
  \label{eq:slip-angle}
\end{equation}

\textbf{Physical Meaning}: Slip angle is the angle between the tire's heading and its velocity direction. It quantifies how much the tire is "crabbing" sideways, a key input to the lateral force model.

\textbf{Source Code}: \coderef{physics.js}{641--643}

\begin{lstlisting}[caption={Slip angle calculation}]
const minSlipDenom = 0.5;  // m/s — avoid singularity at zero speed
const slipDenom = Math.max(Math.abs(vLong), minSlipDenom);
const slipAngle = Math.atan2(vLat, slipDenom);
\end{lstlisting}

\textbf{Parameters}:
\begin{itemize}
  \item $v_{\text{lateral}}$: Lateral velocity component (m/s)
  \item $v_{\text{longitudinal}}$: Longitudinal velocity component (m/s)
  \item $\epsilon = 0.5$ m/s: Minimum denominator to prevent division by zero
\end{itemize}

\subsection{Low-Speed Lateral Force Fade}

\begin{equation}
  F_{\text{lat, faded}} = F_{\text{lat}} \cdot \clampz\left(\frac{v_{\text{total}}}{v_{\text{fade}}}\right)
  \label{eq:low-speed-fade}
\end{equation}

\textbf{Physical Meaning}: Below 2 m/s, lateral forces are scaled down to zero to prevent artificial "shaking" caused by constraint solver micro-impulses being amplified by the nonlinear Pacejka curve.

\textbf{Source Code}: \coderef{physics.js}{653--662}

\textbf{Parameters}:
\begin{itemize}
  \item $v_{\text{fade}} = 2.0$ m/s: Fade threshold
  \item $v_{\text{total}}$: Total speed magnitude
\end{itemize}

\textbf{Notes}: This is purely a cosmetic fix at real driving speeds (onset at 7 km/h).

\newpage
